\section{Connecting People and Things: Shared Worlds, Touch and Telepresence}\label{sec:connect}

Once several people or machines share one world, delivering media is not enough. The network must keep endpoints consistent and responsive to each other: a gesture made at one site must arrive soon enough to be answered, and an object moved by one user must not rest elsewhere for another. Three kinds of loop run through such a session, and different things bound them (\cref{fig:connect-arch,tab:connect-loops}). A control loop, such as bilateral haptic teleoperation, closes a feedback system across the network, so delay threatens stability before it threatens experience. Interactive loops carry shared state and remote people; perception bounds them, because delay matters when users notice it or when it changes how they act together. Freshness loops keep a digital twin in step with a physical asset; the decision the viewer must take bounds them. The section is organised by what is shared: state, people, touch, and machines. Haptic and holographic representations belong to \cref{sec:encode}, the radio and transport beneath these flows to \cref{sec:provision}, and group experience to \cref{sec:assure}.

\begin{figure}[t]
\centering
\begin{tikzpicture}[font=\sffamily\scriptsize, x=1cm, y=1cm,
  >={Stealth[length=1.8mm]},
  box/.style={draw=connect!70!black, fill=connect!10, rounded corners=2pt, align=center,
              line width=0.5pt, inner sep=2pt, text=black},
  ebox/.style={box, fill=connect!18, minimum width=3.6cm, text width=3.4cm, minimum height=0.95cm},
  lbox/.style={box, minimum width=2.8cm, text width=2.6cm, minimum height=0.95cm},
  ctrl/.style={<->, line width=1.3pt, draw=connect!75!black},
  inter/.style={line width=0.65pt, draw=black!75},
  fresh/.style={line width=0.65pt, draw=black!65, dash pattern=on 3pt off 2pt},
  lab/.style={font=\sffamily\scriptsize, align=center, text=black, inner sep=1.5pt},
  num/.style={lab, text=black!70},
  tier/.style={font=\sffamily\scriptsize\bfseries, text=black!80}]
% tier headers
\node[tier] at (1.4,0.85) {Site A: co-located};
\node[tier] at (7.7,0.85) {Edge server (MEC)};
\node[tier] at (14.8,0.85) {Remote users \& robot};
% edge tier background
\begin{scope}[on background layer]
  \fill[connect!5, rounded corners=3pt] (5.75,0.6) rectangle (9.65,-3.3);
  \draw[connect!55!black, line width=0.4pt, dash pattern=on 2pt off 1.5pt, rounded corners=3pt]
    (5.75,0.6) rectangle (9.65,-3.3);
\end{scope}
% left column
\node[lbox] (cap) at (1.4,0) {Capture rig\\RGB-D cameras};
\node[lbox] (grp) at (1.4,-1.35) {Co-located AR group\\clients A, B on Wi-Fi};
\node[num, below=2pt of grp] (anc) {shared anchors: map merge\\${\leq}$200~ms, transfer 12.7--87~s};
\node[lbox] (op) at (1.4,-3.95) {Operator\\haptic master device};
% edge column
\node[ebox] (fus) at (7.7,0) {Fusion \& render\\views $\rightarrow$ 3D $\rightarrow$ viewport};
\node[ebox] (st) at (7.7,-1.35) {State / relay server\\interest management,\\dead reckoning};
\node[ebox] (tw) at (7.7,-2.75) {Digital twin\\age-of-information updates};
% right column
\node[lbox, minimum height=2.3cm] (rem) at (14.8,-0.675) {Remote VR user\\HMD on 5G\\[3pt]sees avatars,\\holograms, twin state};
\node[lbox, minimum height=2.05cm] (rob) at (14.8,-3.35) {Remote robot\\slave arm, sensors};
% holoportation path (interactive)
\draw[inter, ->] (cap.east) -- node[lab, above] {RGB-D streams} node[num, below] {uplink 3--100~Mbps} (fus.west);
\draw[inter, ->] (fus.east) -- node[lab, above] {rendered view} node[num, below] {E2E $\approx$70~ms edge decode,\\170~ms phone decode} (fus.east -| rem.west);
% shared state (interactive)
\draw[inter, <->] (grp.east) -- node[lab, above] {pose, avatar, objects} node[num, below] {31--83~kbps per user} (st.west);
\draw[inter, <->] (st.east) -- node[lab, above] {forward-all replication} node[num, below] {2 users: 101.7--239.1~ms} (st.east -| rem.west);
% freshness loop
\draw[fresh, ->] (tw.east -| rob.west) -- node[lab, above] {telemetry} node[num, below] {29.63--1,015~ms} (tw.east);
\draw[fresh, ->] (tw.north) -- (st.south);
\node[lab, anchor=west] at ($(tw.north)!0.5!(st.south)+(0.05,0)$) {twin state};
% control loop
\draw[ctrl] (op.east) -- node[lab, above, pos=0.5] {force and position samples at 1~kHz}
  node[num, below, pos=0.5] {1~ms target; QoE held to 200~ms with TDPA/MMT switching} (op.east -| rob.west);
% legend
\begin{scope}[shift={(0.6,-5.0)}]
  \draw[ctrl] (0,0) -- (0.8,0); \node[lab, anchor=west] at (0.9,0) {control loop (stability-bound)};
  \draw[inter, ->] (5.3,0) -- (6.1,0); \node[lab, anchor=west] at (6.2,0) {interactive loop (perception-bound)};
  \draw[fresh, ->] (10.9,0) -- (11.7,0); \node[lab, anchor=west] at (11.8,0) {freshness loop (decision-bound)};
\end{scope}
\end{tikzpicture}
\caption{Architecture of a multi-party XR session: a co-located group and a remote user share state through an edge relay, a capture rig feeds edge fusion for holoportation, and an operator drives a remote robot whose digital twin updates every viewer. The loops differ in what bounds them (line style), so no single latency target serves the session. Numbers are measured values cited in the text: avatar rates and replication latency \cite{W4307020344}, holoportation rates and latency \cite{W4312662428,W4307263396}, anchors \cite{W4311118480,W4411499419}, twin updates \cite{W7213228059,W7168263637}, haptic rate and tolerance \cite{W2912555763,W2216406802,W4200154132}.}
\label{fig:connect-arch}
\end{figure}

\begin{table}[t]
\centering\small
\caption{The three loop types of a shared XR session. Measured tolerances come from the studies cited in the text; they show that each loop fails in a different way and needs its own network target.}
\label{tab:connect-loops}
\begin{tabularx}{\linewidth}{@{}>{\raggedright\arraybackslash}p{2.05cm}>{\raggedright\arraybackslash}X>{\raggedright\arraybackslash}p{2.0cm}>{\raggedright\arraybackslash}X>{\raggedright\arraybackslash}X@{}}
\toprule
Loop & What it carries & Bound & Delay measured or tolerated & Failure mode \\
\midrule
Control (haptic teleoperation) & Force and position samples at 1~kHz & Stability of the bilateral controller & 1~ms target; stable to 200~ms with TDPA/MMT switching \cite{W4200154132}, model-mediated to 1~s round trip \cite{W2811040154} & Oscillation or unsafe contact force \\
Interactive: shared state & Pose, avatar and object updates, tens of kbps per user & Perception of consistency and co-presence & 101.7--239.1~ms on deployed platforms \cite{W4307020344}; impairment noticed, task time unchanged \cite{W4392600587} & Divergent views, conflicting edits, unfair interaction \\
Interactive: remote people & Avatar parameters to volumetric video, 0.2--100~Mbps & Perception of the person & 70--170~ms end to end \cite{W4312662428}; avatar latency noticed at about 320~ms \cite{W4312899133} & Degraded, frozen or vanished person \\
Freshness (digital twin) & Telemetry and state of a physical asset & The viewer's decision & 29.63--1,015~ms, all reported adequate \cite{W7213228059,W7168263637} & Decisions taken on stale state \\
\bottomrule
\end{tabularx}
\end{table}

\subsection{Shared state: who holds the truth, and who sees what}

Every shared world must decide where the authoritative copy of its state lives and how much of it each client receives. Three families answer differently. \emph{Authoritative relays} order all updates at a server and forward them. Deployed social VR uses the simplest form: five commercial platforms forward every avatar update unprocessed, so throughput, latency and client load grow roughly linearly with users, and two-user end-to-end latency was 101.7--239.1~ms \cite{W4307020344}. A relay gives one ordering of events but puts itself in every path; routing through a cloud middleware raised round-trip time from 0.04~s to 0.6~s \cite{W4386499017}. \emph{Replicated, optimistic state} lets each client apply its own actions at once and reconcile later. Croquet-style virtual time \cite{W3174543121} and decoupled edge physics serving over 100 users \cite{W4404050807} belong here. They buy responsiveness and pay in conflicts: two users actuating one digital twin without a coordination policy conflicted in 33--56\% of time windows \cite{W7168268649}. \emph{Peer-to-peer and information-centric sharing} removes the relay. Named-data synchronisation \cite{W4404102246}, device-to-device exchange for web AR \cite{W3117560819} and co-located multicast, which rebuilt scenes several times faster than unicast TCP for groups of 7 and 16 \cite{W7147356863}, take this route. No study has compared the three families on the same workload.

Which inconsistencies users notice is assumed, not measured; consistency models appear only as concepts \cite{W3115046446,W2951580915}. Augmented Virtual Teleportation deliberately left its 1.2~s live video unsynchronised with avatar state to stay responsive \cite{W3005948101}. Yet synchronisation faults were the most common symptom in 2,649 bug reports from multi-user XR software \cite{W7166703064}.

Scaling fails at replication before it fails at link capacity. In a controlled study, update intervals grew from 18~ms with 2 clients to 158~ms with 125, and perceived fluidity collapsed beyond 25 avatars \cite{W2922194619}. Networked games use two remedies. Interest management sends each client only the state near or relevant to it. Dead reckoning sends a motion model and updates it only when the receiver's extrapolation error exceeds a bound. XR synchronisation has barely adopted either, and neither has been tested with XR users; lossy compression of joint rotations, a partial step, cut delay from over 300~ms to under 25~ms at 100~Hz \cite{W3008045703}. Large-scale claims rest on dummy clients \cite{W7130552889}.

Co-located AR adds a shared coordinate frame, which becomes the bottleneck. SPAR cut spatial inconsistency by up to 60\% by choosing what to communicate \cite{W3110126341}, and edge map merging finished within 200~ms \cite{W4311118480}, yet anchor transfers took 12.7--87.18~s \cite{W4411499419} and multi-user AR over public LTE had a 12.5~s median end-to-end latency \cite{W3047439371}.

Fairness between users is the most contested question. Edge schedulers minimise inter-player delay \cite{W4285149407,W3160366001}, and predictive buffering cut inter-client skew by 76\% for 100 clients in a controlled testbed \cite{W7168249167}. They assume that equal delay means equal experience. User evidence points elsewhere: dyads perceived network impairment but completed tasks no slower \cite{W4392600587}. The largest controlled study used 30 pairs \cite{W7163944086}, and a review of 56 studies found that almost all optimise latency for an individual, not a group \cite{W7154321690}. Group quality of experience (QoE), treated in \cref{sec:assure}, is unmeasured beyond pairs.

\subsection{People: a ladder of fidelity}

A remote person can be sent as a few joint angles or as a live volumetric capture; each rung up buys realism with bits, latency and capture hardware. \emph{Pose and parametric avatars} are cheapest: 24~B per update for a head pose and 216~B for a 17-joint skeleton \cite{W2922194619}, or 31--83~kbps per user on VRChat, Rec Room and Hubs \cite{W4307020344}. Expressiveness multiplies the rate; Horizon Worlds exceeded 4.5~Mbps downstream at 15 users \cite{W4307020344}. \emph{Semantic, model-driven avatars} send features that drive a pre-built model of the person. Mon3tr animated a Gaussian avatar with under 0.2~Mbps at about 80~ms \cite{W7123904049}, and MagicStream cut bandwidth up to 1,195-fold by sending keypoint and colour semantics \cite{W4404034698}. Apple's spatial personas use about 0.7~Mbps without rate adaptation \cite{W4403981191}. This family assumes an enrolled model and a stable feature stream; its robustness to loss is unsettled.

\emph{Volumetric telepresence} captures, uplinks, fuses and renders live geometry, and each stage has its cost. Raw point clouds of one person need about 300~Mbps; tiled user-adaptive streaming brought this to 14~Mbps \cite{W4304091566}. An operator trial streamed 5--100~Mbps, with end-to-end latency near 170~ms when phones decoded and 70~ms when the edge did \cite{W4312662428}; depth packed into H.264 needed 3--30~Mbps \cite{W4307263396}. Edge multipoint control units cut client load as parties grow \cite{W4387969165}. At the top, \emph{light-field booths} such as Project Starline \cite{W4200107448} need dedicated rigs; a full-duplex light-field link used 173~Mbps on a LAN \cite{W3176472668}.

Reusing the 2D real-time stack works: a WebRTC selective forwarding unit for point clouds adapted faster than TCP \cite{W7151958829}, and L4S converged on bandwidth 45\% faster than Google congestion control without frame loss \cite{W4408870658}. Two assumptions go unexamined. Calls are treated as downlink problems, yet only two studies measured the holographic uplink \cite{W4210904516,W4416047836}, and the operator trial attached its camera by Ethernet \cite{W4312662428}. Semantic schemes report savings in PSNR, never against geometry at equal bitrate with users. Whether 3D beats 2D is contested. Project Starline \cite{W4200107448} and stroke consultations with 30 outpatients \cite{W4400580381} favour 3D, while a floating-head prototype \cite{W7127277161} and a VR2Gather--Zoom comparison \cite{W4415708482} rated 3D below 2D video. All held the network fixed, although humanoid avatars lost more QoE under network distortion than cartoon ones \cite{W4415821547}.

\subsection{Touch: closing a control loop across the network}

Haptic teleoperation is the only loop here that closes a feedback system, so delay can destabilise it. The field set its requirement early at 1~ms end-to-end latency with 99.999\% reliability \cite{W2216406802}, and IEEE 1918.1 fixed a master--network--slave reference architecture \cite{W2911218692}. The requirement is costly at the radio: under finite-blocklength coding, transmit power becomes unbounded once the queueing-delay bound falls below the channel coherence time \cite{W2963776947}. Control theory offers three families that relax it. \emph{Wave variables} encode force and velocity so that the channel stays passive, but they handle only constant delay. \emph{Time-domain passivity} (TDPA) monitors the energy the network injects and dissipates it; with perceptual coding it stayed stable under time-varying delay of 100~$\pm$~30~ms while sending up to 80\% fewer packets. \emph{Model-mediated teleoperation} (MMT) renders force from a local model of the remote environment and sends only model updates; it cut data by about 90\% and stayed stable at round-trip delays up to 1~s \cite{W2811040154}. MMT assumes an environment that changes slowly enough to be modelled; TDPA suits fast-changing contact at low delay \cite{W2963793563}.

Perceptual coding is orthogonal to these families. Deadband coders send a sample only when the change would be felt, using Weber fractions of 7--15\% for force, and reduce packet rates by 80--90\% \cite{W2912555763}; codec details belong to \cref{sec:encode}. Prediction carries the same idea to the network: co-designing receiver-side prediction with radio resources saved up to 77.80\% bandwidth \cite{W4321488300}, and octree-based MMT raised force accuracy by 44\% with fewer than half the model updates \cite{W4401416073}. Both were evaluated open-loop on recorded or synthetic trajectories, not with a live network, a stable controller and users.

Measured tolerance is far looser and depends on the controller. Switching between TDPA and MMT held QoE at delays up to 200~ms \cite{W4200154132}, and users noticed short haptic delays but compensated \cite{W3162230828}. The 1~ms figure assumes a controller without delay compensation; it is not a measured limit of human touch. Transport results split by medium: a time-sensitive networking testbed kept high-priority flows at dedicated-wire quality under 200\% over-saturation \cite{W4387917733}, whereas Wi-Fi needs MAC changes that reached a 3.72~ms worst case for Wi-Fi~7 only in simulation \cite{W4205674533}. Haptic streams rarely travel with their video, and cross-modal work questions their usual priority: deliberate haptic loss with visual restoration raised opinion scores by more than 80\% at 30\% loss \cite{W4401507901}. Multi-user haptic studies stop at six pairs \cite{W4413393831}.

\subsection{Machines and places: keeping a twin fresh}

A digital twin reverses telepresence: the network must keep a replica of a physical asset fresh enough for the viewer's decision. The natural metric is age of information (AoI), the time since the newest received update was generated. Two families use it differently. \emph{Freshness optimisers} allocate radio and edge resources to minimise AoI or synchronisation time. They optimised AoI and reliability for short industrial packets \cite{W4366251214}, and optimal-transport association cut synchronisation time from 31.23--51.23~ms to 19.11--25.17~ms \cite{W4387870504}. They assume their workloads, such as 10~Mb per update \cite{W4387870504}, and rarely leave simulation; a task-oriented design on a real robot arm, which halved the required packet rate, is an exception \cite{W4390044355}. \emph{Measured pipelines} report what deployed stacks deliver. A robot vehicle mirrored on a headset updated in 29.63~ms over a local network \cite{W7213228059}, but middleware dominated elsewhere: 416~ms per transmission for a robot-arm VR twin \cite{W4417338851} and 1,015~ms end to end for a cobot through a commercial IoT platform \cite{W7168263637}. Each figure was called adequate; none was set against a task requirement.

Sending state rather than pixels is a consistent lever. Pose streaming kept remote work usable under added latency where video did not \cite{W4312316324}, a persistent twin with transient object updates allowed teleoperation at kilobits per second \cite{W4405784937}, and server push used three orders of magnitude less bandwidth than polling for a 5,005-building city twin \cite{W7172377742}. Once users act through the twin, the freshness loop becomes a shared-state problem, with the conflicts seen above \cite{W7168268649}. Freshness metrics have multiplied, yet none has been tested against what a viewer perceives or decides.

\begin{openproblems}
\item \textbf{What makes a group experience fair?} Skew schedulers assume equal delay means equal experience, but controlled studies stop at pairs \cite{W7163944086,W4392600587}. Groups of two to four under per-user impairments, scored on a shared task and perceived fairness, would give schedulers a validated target.
\item \textbf{Can interest management and dead reckoning scale social XR?} Platforms forward every update and scale linearly \cite{W4307020344}. Area-of-interest filtering with error-bounded extrapolation on an open platform, with a perceptual study of avatar error, would show how far replication cost can fall.
\item \textbf{How much delay does touch tolerate in XR tasks?} The 1~ms target persists while users tolerated 200~ms with adaptive control \cite{W4200154132}. Factorial delay, loss and jitter studies across TDPA, MMT and closed-loop predictors would yield tolerance curves for provisioning.
\item \textbf{How should haptic flows share a link with XR video?} Haptics and remote-rendered frames are scheduled apart. Joint scheduling on Wi-Fi or 5G, measuring haptic tail delay, motion-to-photon latency and inter-stream skew, would answer it.
\item \textbf{Which rung of the fidelity ladder serves a conversation?} Semantic avatars were never compared with geometry at equal bitrate, and 3D-versus-2D studies hold the network fixed. Two-way calls varying delay, loss and uplink rate against a matched 2D control are needed.
\item \textbf{How stale may a twin be?} AoI is optimised, yet no study links staleness to detection or control errors \cite{W4387870504,W7168263637}. Psychometric staleness curves for monitoring and actuation tasks would turn freshness budgets into requirements.
\end{openproblems}
